\begin{textsample}{POS Dim 1 – ms – Score 23.00 – ms/ms\_inf\_e\_ef\_32.txt}  \label{ex:f1_pos_210}
8.
ENSINO FUNDAMENTAL > A maior riqueza do \textbf{homem} é sua incompletude. > Nesse ponto sou abastado. > Palavras \textbf{que} me aceitam > como sou — eu não aceito. > Não aguento ser apenas > \textbf{um} \textbf{sujeito} \textbf{que} abre portas, \textbf{que} puxa > válvulas, \textbf{que} olha o relógio, \textbf{que} compra pão > às 6 da tarde, \textbf{que} vai > lá fora, \textbf{que} aponta lápis, > \textbf{que} \textbf{vê} a uva etc. etc. > Perdoai.
Mas eu > \textbf{preciso} ser Outros. > Eu penso \textbf{renovar} o \textbf{homem} > usando borboletas. > > Manoel de Barros ( 1998 ) O Ensino Fundamental, etapa intermediária da Educação Básica, com duração de 9 anos, tem como \textbf{fundamento} o acesso ao conhecimento \textbf{historicamente} elaborado e a elementos culturais \textbf{que} asseguram o desenvolvimento humano tão necessário para o convívio em sociedade.
No entanto, cabe mencionar \textbf{que} durante esta etapa, os \textbf{sujeitos} transformam -se, o \textbf{que} quer \textbf{dizer}, \textbf{que} mudanças nos aspectos físicos, cognitivos, afetivos, sociais e emocionais ocorrem nas crianças e nos adolescentes.
Pode -se traduzir esta situação ao mencionar \textbf{que} a criança acessa o Ensino Fundamental aos 6 anos e finaliza essa etapa aos 14 anos, ou seja, durante o processo de escolarização no Ensino Fundamental o \textbf{sujeito} tem elementos da infância e da adolescência. \textbf{Dito} isso, o Currículo deve minimizar os impactos dessas fases e transformações, considerando os ritmos dos \textbf{sujeitos} e possibilitando estratégias de aprendizagem \textbf{que} não demonstrem rupturas entre as fases do Ensino Fundamental ( anos iniciais e anos finais ).
O Currículo dos anos iniciais do Ensino Fundamental reconhece a necessária articulação com as experiências vivenciadas pela criança na Educação Infantil e preza pelas situações lúdicas de aprendizagem.
Assim, as estratégias de aprendizagem devem \textbf{sistematizar} as experiências das crianças com vistas à ampliação dos conhecimentos e das relações \textbf{que} estão estabelecendo consigo mesmas, com os outros e com o mundo.
Para tal, o Parecer das Diretrizes Curriculares Nacionais indica \textbf{que} : > A escola deve adotar formas de trabalho \textbf{que} proporcionem maior mobilidade às crianças na sala de aula, explorar com elas mais intensamente as diversas linguagens artísticas, a começar pela literatura, utilizar mais materiais \textbf{que} proporcionem aos estudantes oportunidade de racionar manuseando -os, explorando as suas características e propriedades, ao mesmo tempo em \textbf{que} passa a \textbf{sistematizar} mais os conhecimentos escolares ( BRASIL, p. 121, 2013 ).
Além dessas características \textbf{que} visam à redução dos impactos na criança da ruptura entre a pré-escola e o Ensino Fundamental, há \textbf{que} se considerar \textbf{que} ela está envolvida em processos de alfabetização e letramento e estes não podem ser interrompidos ao final do primeiro ano do Ensino Fundamental.
Aos \textbf{sujeitos} pertencentes aos anos iniciais do Ensino Fundamental possibilita -se a ampliação da interação com os espaços ; as relações com as múltiplas linguagens ; as relações sociais, a afirmação da identidade e a valorização das diferenças.
Dessa forma, “ a criança desenvolve a capacidade de representação, indispensável para a aprendizagem da leitura, dos conceitos matemáticos básicos e para a compreensão da realidade \textbf{que} a cerca, conhecimentos \textbf{que} se postulam para esse período da escolarização ” ( BRASIL, p. 110, 2013 ).
Ainda, por meio do desenvolvimento da linguagem, as crianças reconstroem, descrevem e planejam suas ações com o auxílio da memória.
Esse desenvolvimento, no Currículo de Mato Grosso do Sul, aponta para os processos de aprendizagem \textbf{que} envolvem a aquisição da leitura e da escrita vinculada aos seus usos sociais ; à \textbf{percepção} e representação das experiências pela oralidade ; pelos signos matemáticos ; pelos registros artísticos ; pela cultura digital e pela ciência.
Nesse sentido, a educação integral, prevista neste Currículo, aponta para processos educativos \textbf{que} favoreçam nas crianças o pensamento criativo, lógico e crítico durante o uso das múltiplas linguagens.
O foco da alfabetização ocupa espaço \textbf{central} nos dois primeiros anos do Ensino Fundamental e tem como objetivo garantir às crianças a apropriação do sistema de escrita alfabética e o envolvimento em práticas de letramentos.
Assim, todos os conhecimentos mobilizados pelos componentes curriculares dos anos iniciais do Ensino Fundamental contribuem para a consolidação desse objetivo.
Conforme o Parecer das Diretrizes Curriculares Nacionais para o Ensino Fundamental de 9 anos : > [... ] desde os 6 ( seis ) anos de idade, os conteúdos dos demais componentes curriculares devem também ser trabalhados.
São eles \textbf{que}, ao descortinarem às crianças o conhecimento do mundo por meio de novos olhares, lhes oferecem oportunidades de exercitar a leitura e a escrita de \textbf{um} modo mais significativo ( BRASIL, p. 22, 2010 ).
Durante os anos iniciais do Ensino Fundamental, o Currículo de Referência de Mato Grosso do Sul organiza -se para \textbf{que} sejam consolidadas as aprendizagens anteriores e para a “ ampliação das práticas de linguagem e da experiência estética e intercultural das crianças ” ( BRASIL, p. 57, 2017 ).
Desse modo, a autonomia intelectual, pressuposto da educação integral, amplia -se e, consequentemente, \textbf{renovam} -se os interesses pela vida social no tocante “ às relações dos sujeitos entre si, com a natureza, com a história, com a cultura, com as tecnologias e com o ambiente ” ( BRASIL, p. 57, 2017 ).
Os \textbf{sujeitos} dos anos finais do Ensino Fundamental correspondem a \textbf{uma} faixa de transição entre a infância e a adolescência de extremas mudanças biológicas, sociais, psicológicas e emocionais.
Segundo o Parecer CNE/CEB n. 11/2010 “ os adolescentes, nesse período da vida, modificam as relações sociais e os laços afetivos, intensificando suas relações com os pares de idade e as aprendizagens referentes à \textbf{sexualidade} e às relações de gênero, acelerando o processo de ruptura com a infância na tentativa de construir valores próprios ” ( BRASIL, p. 09, 2010 ).
Com isso, as possibilidades intelectuais são ampliadas o \textbf{que} ocasiona a realização de \textbf{raciocínios} mais abstratos.
Sendo assim, os desafios são mais complexos nos anos finais do Ensino Fundamental, especialmente porque os conhecimentos adquiridos até o momento serão aprofundados.
A ação pedagógica requer a retomada e a ressignificação de aprendizagens, a fim de ampliar o conhecimento \textbf{historicamente} acumulado dos adolescentes. \textbf{Ademais}, os processos educativos devem considerar as mudanças próprias da adolescência, reconhecendo \textbf{que} esses \textbf{sujeitos} estão em desenvolvimento, ou seja, promovendo práticas escolares \textbf{dialógicas} \textbf{que} contemplem as diversas formas de inserção social.
Sobre isso, as Diretrizes Curriculares Nacionais para o Ensino Fundamental indicam \textbf{que} : > Entre os adolescentes de muitas escolas, é frequente observar \textbf{forte} adesão aos padrões de comportamento dos jovens da mesma idade, o \textbf{que} é evidenciado pela forma de se vestir e também pela linguagem utilizada por eles.
Isso requer dos educadores maior disposição para entender e dialogar com as formas próprias de expressão das culturas juvenis [... ] ( BRASIL, p. 110, 2013 ).
Outra questão relevante refere -se às transformações sociais originadas pela cultura digital.
Sabe -se \textbf{que} os jovens interagem socialmente em rede, estabelecem novas formas de interação multimidiática e multimodal, e, por \textbf{conseguinte}, são protagonistas dessa cultura.
Entretanto, a cultura digital recorre às emoções, ao imediatismo, às informações superficiais, ao uso de imagens e às formas de expressão abreviadas.
Nessa perspectiva, a escola, compromissada com a formação integral do \textbf{sujeito}, pode incorporar as novas linguagens advindas do desenvolvimento das Tecnologias Digitais de Informação e Comunicação, assim como oportunizar a reflexão e análise de conteúdos midiáticos, por meio de \textbf{uma} educação consciente para/na cultura digital.
A educação fundamenta -se nos direitos humanos e em princípios democráticos, o \textbf{que} \textbf{revela} a preocupação e o combate a qualquer forma de violência propagada nesta sociedade.
Dá -se atenção para essa problemática nos anos finais do Ensino Fundamental devido à fragilidade dos \textbf{sujeitos} \textbf{que}, quando em situação de violência, têm a convivência cotidiana e a aprendizagem comprometidas \textbf{que}, em última \textbf{instância}, \textbf{revela} -se em fracasso escolar e autodestruição. \textbf{Logo}, o Currículo de Referência de Mato Grosso do Sul \textbf{fomenta} o diálogo entre as diferentes culturas existentes na escola, a compreensão da não uniformidade dos \textbf{sujeitos}, ao reconhecimento da diversidade de experiências e vivências porque entende a escola como espaço para formação da cidadania participativa.
Para esse fim, o Currículo dos anos finais do Ensino Fundamental contribuirá na elaboração do projeto de vida dos adolescentes e jovens e, consequentemente, auxiliará nas suas decisões em relação ao seu futuro e na continuidade dos estudos no Ensino Médio.

% matched lemmas: ademais, central, conseguinte, dialógico, dizer, fomentar, forte, fundamento, historicamente, homem, instância, logo, percepção, preciso, que, raciocínio, renovar, revelar, sexualidade, sistematizar, sujeito, um, ver
\end{textsample}
